% Created 2026-07-08 Wed 23:50
% Intended LaTeX compiler: pdflatex
\documentclass{article}
\usepackage[utf8]{inputenc}
\usepackage[T1]{fontenc}
\usepackage{graphicx}
\usepackage{longtable}
\usepackage{wrapfig}
\usepackage{rotating}
\usepackage[normalem]{ulem}
\usepackage{amsmath}
\usepackage{amssymb}
\usepackage{capt-of}
\usepackage{hyperref}

\usepackage{fancyhdr}
\usepackage[top=25mm, left=25mm, right=25mm]{geometry}
\usepackage{longtable}
\fancyhead[R]{}
\setlength\headheight{43.0pt}
\usepackage[T1]{fontenc}
\usepackage[utf8]{inputenc}
\usepackage[spanish, ]{babel}
\usepackage[backend=biber]{biblatex}

\usepackage{dirtree}
\author{Jami Tonato Mateo Israel}
\date{2026-07-08}
\title{Taller Script Python}
\hypersetup{
 pdfauthor={Jami Tonato Mateo Israel},
 pdftitle={Taller Script Python},
 pdfkeywords={},
 pdfsubject={},
 pdfcreator={Emacs 27.1 (Org mode 9.7.5)}, 
 pdflang={Español}}
\begin{document}

\maketitle
\fancyhead[C]{\includegraphics[scale=0.05]{../images/logoEPN.jpg}\\
ESCUELA POLITÉCNICA NACIONAL\\FACULTAD DE INGENIERÍA DE SISTEMAS\\
ARQUITECTURA DE COMPUTADORES}
\thispagestyle{fancy}


\section{Objetivos}
\label{sec:orgf4f2344}
\begin{itemize}
\item Desarrollar un script de Python con paradigma estructurado
\end{itemize}

\section{Problema}
\label{sec:org5d7bae2}
Usted es el científico de datos para un sistema de lanzamiento de
cohetes. Para esto, se le solicita que determine gráficamente el
alcance del cohete según al menos tres ángulos de entrada, un tiempo
máximo de vuelo \texttt{t}. Se considera que existe una fuerza de resistencia
al movimiento debido al viento y representada mediante un coeficiente
\texttt{b}. Escribir un programa que permita obtener el alcance del cohete
gráficamente, mediante la resolución de ecuaciones diferenciales
ordinarias. Los parámetros como Velocidad inicial \texttt{V}, coeficiente de
resistencia \texttt{B}, tiempo máximo \texttt{tmax} y ángulos de lanzamiento se
ingresan desde la línea de comando. Siga las instrucciones
proporcionadas en  la Sección \ref{sec:orgeb4a700} para
completar el taller.

\section{Instrucciones}
\label{sec:orgeb4a700}

\begin{enumerate}
\item Obtener las expresiones matemáticas para las ecuaciones
diferenciales ordinarias. Vea la Sección \ref{sec:org8f538c6}
\item Investigar cómo funciona la librería \texttt{solve\_ivp} del paquete
\texttt{scipy.integrate}
\item Investigar cómo funciona la librearía \texttt{argparse} para controlar el
ingreso de valores y opciones desde la línea de comandos.
\item Cree una carpeta con los archivos del proyecto denominada \texttt{cohete}.
\item Dentro de la carpeta cree un archivo \texttt{main.py} con el código
proporcionado en la Sección \ref{sec:orgff1aa88}
\item Reorganizar el código proporcionado de ejemplo para conformar un
proyecto de código estructurado. Las funciones estarán dentro de
una carpeta \texttt{utils} en el archivo \texttt{projectile\_functions.py}. El
código principal se coloca en el archivo \texttt{main.py} conforme al
siguiente esquema de archivos.
\begin{enumerate}
\item Con YASnippet habilitado escriba \texttt{imp} seguido de \texttt{C-TAB} para
escribir \texttt{import}
\item Escriba \texttt{main} seguido de la secuencia
de comandos \texttt{C-TAB}. Esto genera una función llamada \texttt{main}
\item escriba \texttt{ifm} seguido de la secuencia \texttt{C-TAB} para generar la
condición de ejecución del programa Python como programa
principal
\item Use \texttt{fd} seguido de \texttt{C-TAB} para escribir una función con
\emph{docstrings} y use \texttt{r} seguido de \texttt{C-TAB} para escribir \texttt{return}
\item Utilice \texttt{pars} \texttt{C-TAB} para crear un parser y \texttt{args} \texttt{C-TAB}
para crear un argumento de entrada del programa. Importe la
librería \texttt{argparse}
\end{enumerate}
\end{enumerate}

\dirtree{%
 .1 /. 
 .2 main.py. 
 .2 utils. 
 .3 \_\_init\_\_.py. 
 .3 projectile\_functions.py. 
 }
\section{Entregables del taller}
\label{sec:org9568a2f}
\begin{enumerate}
\item Archivo .zip con el archivo \texttt{main.py} y la librería desarrollada
\texttt{projectile\_functions.py} y el gráfico obtenido de una carrera.
\item Completar en la sección [[]] de este documento una breve
descripción de lo que hacen las librerías revisadas en este taller.
\item Una vez redactado el punto anterior suba el archivo .ORG y el .PDF
\end{enumerate}
\section{Recursos}
\label{sec:org3c83a61}
Puede consultar más información en el siguiente \href{https://www.youtube.com/watch?v=b9S\_L1AaJNw}{vídeo} y en \href{https://github.com/lukepolson/youtube\_channel/blob/main/Python\%20Metaphysics\%20Series/vid22.ipynb}{GitHub}

\subsection{Desarrollo Matemático}
\label{sec:org8f538c6}


La fuerza de resistencia del viento se puede considerar que es
proporcional al cuadrado de la velocidad y en sentido opuesto a
ésta. De ahí que las fuerzas actuando en el sistema quedarían
determinadas por la ecuación


\begin{equation}
  \label{}
  \vec{F}_{\text{net}} = \vec{F}_g + \vec{F}_f = -mg\hat{y} - b|\vec{v}|\vec{v}
\end{equation}

La velocidad está dada por \(\vec{v} = \dot{x} \hat{i} + \dot{y} \hat{j}\)

En consecuencia la Fuerza neta del sistema queda:

\begin{equation}
  \label{}
  \vec{F}_{\text{net}} = -mg\hat{y} - b\sqrt{\dot{x}^2 + \dot{y}^2}(\dot{x}\hat{i} + \dot{y}\hat{j})
\end{equation}

Que en forma vectorial se puede escribir como

\begin{equation}
  \label{}
  \vec{F}_{\text{net}} = \begin{bmatrix} - b\sqrt{\dot{x}^2 + \dot{y}^2}\dot{x} \\ -mg - b\sqrt{\dot{x}^2 + \dot{y}^2}\dot{y} \end{bmatrix}
\end{equation}

Dado que \(\vec{F}_{\text{net}} = m\vec{a} = m\left< \ddot{x}, \ddot{y} \right>\), se tiene

\begin{equation}
  \label{}
   m \begin{bmatrix}\ddot{x} \\ \ddot{y} \end{bmatrix} =  \begin{bmatrix} - b\sqrt{\dot{x}^2 + \dot{y}^2}\dot{x} \\ -mg - b\sqrt{\dot{x}^2 + \dot{y}^2}\dot{y} \end{bmatrix}
\end{equation}

De donde se obtienen el par de ecuaciones diferenciales

\begin{equation}
  \label{}
  \ddot{x} = - \frac{b}{m}\sqrt{\dot{x}^2 + \dot{y}^2}\dot{x}
\end{equation}

\begin{equation}
  \label{}
  \ddot{y} = -g - \frac{b}{m}\sqrt{\dot{x}^2 + \dot{y}^2}\dot{y}
\end{equation}

Definiendo \(x' = x/g\) and \(y'=y/g\), se tiene entonces

\begin{equation}
  \label{}
   \ddot{x'} = - \frac{bg}{m}\sqrt{\dot{x'}^2 + \dot{y'}^2}\dot{x'}
\end{equation}

\begin{equation}
  \label{}
   \ddot{y'} = -1 - \frac{bg}{m}\sqrt{\dot{x'}^2 + \dot{y'}^2}\dot{y'}
\end{equation}

Sea \(B \equiv bg/m\), entonces considerando el cambio de notación y
eliminado la ' para simplificar, se propone resolver un sistema de
ecuaciones de primer orden ODE siguiente:


\begin{align}
  \label{}
  \dot{x} &= v_x\\
  \dot{v_x}  &= - B\sqrt{\dot{x}^2 + \dot{y}^2}\dot{x}\\
  \dot{y} &= v_y\\
  \dot{v_y} &= - B\sqrt{\dot{x}^2 + \dot{y}^2}\dot{y}
\end{align}



\subsection{Código Inicial Python}
\label{sec:orgff1aa88}


\begin{verbatim}
import numpy as np
from scipy.integrate import solve_ivp
import matplotlib.pyplot  as plt

def dSdt(t,S,B):
    """
    Use fd seguido de C-TAB para generar la funcion con docstrings
    Define el sistema de ecuaciones diferenciales para el movimiento del proyectil
    con resistencia al aire
    Parametros:
        t (float): Tiempo
        S (list): Estado actual [x, vx, y, vy]
        B (float): Coeficiente de resistencia del aire. Varía desde 0.0 a 1.0
    Retorna:
        list: Derivadas [dx/dt, dvx/dt, dy/dt, dvy/dt]
    """
    x, vx, y, vy = S
    return [vx,
            -B*np.sqrt(vx**2+vy**2)*vx,
            vy,
            -1-B*np.sqrt(vx**2+vy**2)*vy]

B = 0
V = 1
t1 = 40*np.pi/180
t2 = 45*np.pi/180
t3 = 50*np.pi/180

sol1 = solve_ivp(dSdt, [0, 2],
                 y0=[0,V*np.cos(t1),0,V*np.sin(t1)],
                 t_eval=np.linspace(0,2,1000),
                 args=(B,)) # atol=1e-7, rtol=1e-4)
sol2 = solve_ivp(dSdt, [0, 2],
                 y0=[0,V*np.cos(t2),0,V*np.sin(t2)],
                 t_eval=np.linspace(0,2,1000),
                 args=(B,)) #atol=1e-7, rtol=1e-4)
sol3 = solve_ivp(dSdt, [0, 2],
                 y0=[0,V*np.cos(t3),0,V*np.sin(t3)],
                 t_eval=np.linspace(0,2,1000),
                 args=(B,)) #atol=1e-7, rtol=1e-4)


plt.plot(sol1.y[0],sol1.y[2], label=r'$\theta_0=40^{\circ}$')
plt.plot(sol2.y[0],sol2.y[2], label=r'$\theta_0=45^{\circ}$')
plt.plot(sol3.y[0],sol3.y[2], label=r'$\theta_0=50^{\circ}$')
plt.ylim(bottom=0)
plt.legend()
plt.xlabel('$x/g$', fontsize=20)
plt.ylabel('$y/g$', fontsize=20)
plt.show()

\end{verbatim}

\section{Descripción librerías usadas}
\label{sec:org9c7a1ba}

\subsection{SOLVE\textsubscript{IVP}}
\label{sec:orgcc10af8}
Esta función es parte del módulo scipy.integrate de la librería
 SciPy. Sirve principalmente para resolver problemas de valor inicial
 en sistemas de ecuaciones diferenciales ordinarias.
 En nuestro proyecto funciona como el motor físico que calcula la
 trayectoria del cohete paso a paso a través del tiempo.
 Para que funcione le pasamos varios parámetros. Primero la función
 dSdt que contiene nuestro sistema de ecuaciones. También le indicamos
 el tiempo de inicio y fin de la simulación con t\textsubscript{span}.
 Las condiciones iniciales de posición y velocidad van en el arreglo
 y0. Usamos t\textsubscript{eval} para decirle que guarde la posición del cohete en
 1000 puntos exactos y así la gráfica salga suave.
 Finalmente con args le pasamos el coeficiente de fricción del aire.

Fuente: SciPy Developers. (2024). scipy.integrate.solve\textsubscript{ivp}. SciPy Documentation. Recuperado de: \url{https://docs.scipy.org/doc/scipy/reference/generated/scipy.integrate.solve\_ivp.html}

\subsection{ARGPARSE}
\label{sec:org9da81ee}

Esta es una librería estándar de Python que se usa para crear
 interfaces de línea de comandos. En el proyecto la usamos para no
 dejar variables quemadas en el código principal.
 Así podemos ingresar la velocidad o los ángulos directamente desde la
 terminal al ejecutar el script sin modificar el código fuente.
 Usamos ArgumentParser para iniciar el contenedor y agregar una
 descripción del programa. Con add\textsubscript{argument} definimos los parámetros
 que vamos a leer.
 Le pusimos el tipo float para que convierta los números
 automáticamente y marcamos los datos como requeridos.
 También usamos nargs="+" para poder ingresar varios ángulos a la vez y que el programa los guarde como una lista. Al final parse\textsubscript{args} recolecta toda esa información para usarla en el código.

Fuente: Python Software Foundation. (2024). argparse. Python
Documentation. Recuperado de:
\url{https://docs.python.org/3/library/argparse.html}

\subsection{NUMPY}
\label{sec:org60eb01e}

NumPy es una librería clave para cálculos matemáticos en Python. En este laboratorio nos ayuda a optimizar los cálculos trigonométricos y a generar arreglos de tiempo sin tener que escribir bucles manuales que harían más lento el programa.

Usamos np.cos y np.sin para descomponer el vector de la velocidad inicial en los ejes horizontales y verticales. También usamos la constante np.pi para transformar los grados que ingresa el usuario a radianes. La función np.sqrt nos sirvió para calcular la raíz cuadrada y sacar la magnitud de la velocidad para aplicar la fuerza de resistencia del aire. Por último usamos np.linspace para generar un arreglo de 1000 números espaciados de forma uniforme y tener una simulación mucho más fluida.

Fuente: NumPy Developers. (2024). NumPy Reference. NumPy Documentation. Recuperado de: \url{https://numpy.org/doc/stable/reference/}
\end{document}
